\section{Attenuation under the principal Gaussian reference}
\label{sec:principal-reference-attenuation}
We apply the polynomial cancellation just proved to the reference
integral in Section~\ref{sec:principal-gaussian-comparison}.
The important feature is the order of constants: the coefficient
of $L^{-d-1}$ is chosen before $L$. Bounds for the nonlinear chart
and the Gaussian can be chosen after $L$, because their additional
costs tend to zero as the entrance coupling increases. Throughout,
$n\ge\max\{2,d+1\}$, the insertion has the fine regularity through
order $n$ used below, and its polynomial moments through degree $d$ vanish.
The coarse weight exponent does not exceed the fine one.

\subsection{The first derivative of the reference mean}
Let $u$ be a coarse logarithmic field and $v_F(u)$ the stationary
relative field in tree coordinates from Section~\ref{sec:local-mean}.
At the identity the full affine row satisfies
\[
                     D_\infty v_F(u)+r'_0(u)=0.
\]
This is \eqref{eq:mean-stationary-zero}, with the same signed
normalization on the two terms. The reference construction uses
the cut row $D_R$ and its cut residual derivative $r'_R$.
The exponential path tail gives, uniformly in the period,
\begin{equation}\label{eq:reference-tangent-tail}
 \|D_Rv_F(u)+r'_R(u)\|_\infty
       \le C(1+R)^a e^{-\sigma R}\|u\|_\infty.
\end{equation}
To see this, subtract the full identity. The row difference acts
on $v_F(u)$, whose sup norm is at most $C_{\rm resp}\|u\|_\infty$.
The guide difference acts on $u$. Each is a sum of omitted path
rows; their weighted absolute tails give the displayed bound,
including the finite path-length power. All inactive coordinates
are omitted on both sides, and tree coordinates of $v_F$ are zero.

Let $A_R=D_R^*D_R$, let $P_R$ be the reference principal-window
inverse, and put $E_R=A_RP_R-I$. The reference mean in active
coordinates is $m(u)=-P_R(u)^*D_R(u)^*r_R(u)$. Since
$r_R(0)=0$, its derivative at zero is
\[
 Dm(0)u=-P_R(0)^*D_R(0)^*r'_R(u).
\]
The elementary normal-equation identity therefore gives
\[
 Dm(0)u-v_F(u)
 =E_R(0)^*v_F(u)
       -P_R(0)^*D_R(0)^*(D_R(0)v_F(u)+r'_R(u)).
\]
The window and row bounds of
Section~\ref{sec:principal-comparison-jets}, together with
\eqref{eq:reference-tangent-tail}, imply
\begin{equation}\label{eq:reference-mean-tangent}
 \|Dm(0)-v_F\|_{\infty\to\infty}
 \le C\{e^{-c\lfloor B/60\rfloor}
                      +(1+R)^a e^{-\sigma R}\}.
\end{equation}
Here and below the restriction to active counting coordinates
and the fixed colour reassembly contribute only fixed factors.
Because $R=\lfloor B/100000\rfloor$ and $B$ grows quadratically
in the logarithm of the coupling, the right side is eventually
at most the coarse small-field threshold $\theta_c$. This proves
the required first-order approximation for the selected windowed
mean, without replacing it by the exact inverse mean.

\subsection{The nonlinear local display}
The tree-straightening potential is a linear site field $\phi(u)$.
Its discrete gradient corrects the tree response to the actual
installed $A_0$ response, denoted $\mathcal A u$:
\[
 \phi(u)(e_-)-\phi(u)(e_+)+v_F(u)_e+\overline u_e
                                             =(\mathcal A u)_e.
\]
This follows by summing the tree potential along a bond: the
interior terms telescope, leaving its endpoint difference. The
root correction is included in the same potential, as in the
response bound of Section~\ref{sec:local-mean}.

For noise $w$ in fine link logarithmic units define the local
logarithmic display
\[
 \Phi(u,w)_e=\log\left(
 e^{\phi(u)(e_-)}
 e^{\widetilde m(u)_e+w_e}e^{\overline u_e}
 e^{-\phi(u)(e_+)}\right),
\]
where $\widetilde m$ is the colour-reassembled mean, zero on tree
links. The logarithm is taken on one common fixed ball. Its
existence and uniform derivatives follow from the fixed-dimensional
group chart, the mean row bounds, the tree-potential bound and the
replica bound. Shrink the input ball so that the four entries of
the logarithm's argument lie in that chart. The derivative of
$\log(e^a e^b e^c e^{-d})$ at zero is $a+b+c-d$.
Consequently, with $T(u,w)=\mathcal A u+w$, there are $Q\ge1$
and $r>0$, uniform over all periods, forests and tags, such that
\begin{equation}\label{eq:reference-display-data}
 \Phi(0)=0,\quad \|D\Phi(0)-T\|\le Q\theta_c,
 \qquad \|D^j\Phi(z)\|\le Q\quad
                  (j\le n,\ \|z\|\le r).
\end{equation}
The constants $Q,r$ may depend on the already chosen block factor.
The tangent error is \eqref{eq:reference-mean-tangent}; the
other bounds follow by the finite product and composition rules
and the separate reference inverse estimates already proved.

Let $a(y)=(u(y),w)$, where $u(y)$ is the fixed-gauge local
logarithm along the ordered coarse variation. Write $s$ for its
common size and derivative scale, so that
\[
 \|a(0)\|\le s,\qquad
 \|D^ja(0)\|\le s^j\quad(1\le j\le n),\qquad
                  \theta_c\le s\le\min\{1,r\}.
\]
The noise is controlled separately from the health condition. In counting
Gaussian coordinates $z$, it is $w=g\,\mathrm{tape}(z)$, up to the fixed
colour isometries. We use a compactly supported $C^n$ cutoff $\Gamma$
with $|\Gamma|\le1$ and
\[
 z\in\operatorname{supp}\Gamma
 \quad\Longrightarrow\quad g\|\mathrm{tape}(z)\|
             \le\varepsilon\theta_f\le s.
\]
The chart bound gives $s=C\theta_c$ times a fixed polynomial in its
read radius, with the local-gauge constant $C$ purchased before $L$.
The nonlinear display constants $Q,r$ are chosen after $L$. Thus the asserted
bound for $a$ holds on the cutoff support, including its boundary.
Only the coarse component varies. Integrating the first two
derivative bounds in \eqref{eq:reference-display-data} gives
\[
 \|\Phi(z)-Tz\|\le Q(\theta_c+\|z\|)\|z\|,
 \qquad \|D\Phi(z)-T\|\le Q(\theta_c+\|z\|).
\]
Apply the partition chain rule to $(\Phi-T)\circ a$. A term
with one block costs at most $2Qs\cdot s^j\le2Qs^2$.
A term with at least two blocks costs at most $Q s^j\le Qs^2$;
linear maps have zero derivatives of those orders. The value is
bounded by $2Qs^2$ too. Thus for a constant $c_n$ depending only
on the derivative order,
\begin{equation}\label{eq:reference-display-remainder}
 \|D^j(\Phi\circ a)(0)-\mathcal A D^ju(0)\|
                    \le c_nQs^2+\|w\|,\qquad j\le n.
\end{equation}
The noise term is needed only at $j=0$ but is harmless at all
orders.

Each $\mathcal A D^ju(0)$ has the graded decomposition of
Section~\ref{sec:polynomial-moment-attenuation}, with cost $M$
chosen before $L$. Put the error in
\eqref{eq:reference-display-remainder} entirely into its last
piece. Since the fine direction width is at least $\theta_f$,
all the display derivatives have a graded decomposition at scale
$L^{-1}$ and cost
\begin{equation}\label{eq:reference-graded-cost}
            M+\frac{c_nQs^2+\|w\|}{\theta_f}L^{d+1}.
\end{equation}
This operation requires no polynomial property of the error.

\subsection{The common healthy segment}
For a short fine support, meaning that its cover number $q$ obeys
$q^3<L$, the whole inspected support lies in the read halo of
one coarse chart. Seed separation, identical local flags, and the
mean-locality argument of
Section~\ref{sec:principal-comparison-jets} identify the displayed
configuration with the physical Gaussian-reference configuration
on that support, after two gauge transformations. One gauge is
the fixed coarse chart gauge; the second is the exponential of
$\phi(u)$. Gauge invariance of the insertion removes both.
The equality holds as a germ in the coarse variation, so derivatives
of the insertion agree. Neither the frozen fine field nor the
Gaussian coordinates are resampled in this identity.

The cutoff restricts to configurations healthy on $X$. We also
need health on the radial segment of the display. Here is the
small elementary estimate that supplies it. If a plaquette has
four logarithmic bonds of norm at most $h\le1$, let $K(t)$ be
the matrix chord
$K(t)=(\operatorname{mat}(\operatorname{plaquette}(t))-I)/c_G$,
where $c_G$ is the fixed colour normalization and the norm is the
ambient matrix norm of the health test. The group product and
this chord map have bounded second derivatives on this fixed ball,
so
\[
 \|K(t)-tK(1)\|\le C_G t(1-t)h^2,
                                      \qquad 0\le t\le1.
\]
Indeed subtract the affine interpolation of $K$ at its endpoints
and integrate its second derivative against the Dirichlet Green
kernel on $[0,1]$; the kernel integral is $t(1-t)/2$.
All four bonds in the touched-plaquette halo are included. The chord
at zero is zero. Hence, if the endpoint is healthy
and $C_Gh^2\le\theta_f$,
\[
                  \|K(t)\|\le t\theta_f+
                       t(1-t)\theta_f\le\theta_f.
\]
Take $h=Qs$. The two required smallness conditions
$Qs\le1$ and $C_G(Qs)^2\le\theta_f$ hold eventually: $s$ is
$\theta_c$ times a fixed polynomial in the read radius, and
$\theta_c^2/\theta_f$ has negative upper exponent. The direction
width bound follows likewise from
$\theta_f(1+6(1+B)L)^p\le1$. These purchases occur after $L$.
We can therefore apply the displaced cancellation bound to every
partition in the derivatives of $f\circ\exp\circ\Phi\circ a$.
It gives the fine-insertion price
\begin{equation}\label{eq:reference-short-pointwise}
 C_{n,d}\tau L^{-d-1}
 \left(M+\frac{c_nQs^2+\|w\|}{\theta_f}L^{d+1}\right)^n.
\end{equation}
The finite number of partitions is included in $C_{n,d}$.

\subsection{Separate local gauges for a winding support}
For a general support $X$, one common logarithm need not exist. We instead
base the insertion chart at the actual received configuration $W_0$.
Keep the frozen fine field, the Gaussian coordinates, and every localization
tag fixed. For each tested fine bond $i$, choose the coarse tree gauge
$g_i$ on the read halo centered at its coarse parent. The local mean and
covariance agreement proved earlier identifies this local calculation with
the same received configuration. The gauges may differ with $i$.

Let $b_i(y)$ be the coarse logarithm in this gauge along the coarse
variation, and let $w_i$ be the adjoint transform of the same noise. Use the unstraightened tree display
\[
 S_i(y)=e^{\widetilde m(b_i(y))_i+(w_i)_i}
                           e^{\overline {b_i(y)}_i}.
\]
Its linear response is the tree response
$\mathcal W b=v_F(b)+\overline b$, whose weighted bound is $M$.
Take $\log(S_i(y)S_i(0)^{-1})$ and return it by the fixed adjoint
$g_i^{-1}$ at the initial endpoint. The variable straightening
potential $\phi(b_i(y))$ is absent from this relative display. Denote the result by $Z_i(y)$;
set it to zero on untested bonds. The relative logarithm is defined
near $y=0$ because both arguments lie in the fixed chart. It satisfies
\[
 Z(0)=0,\qquad
 W(y)_i=\exp(Z_i(y))(W_0)_i
                 \quad\hbox{on every tested bond}.
\]
This is an identity of germs: the relative exponential cancels the
base display in that bond's own gauge. Finitely many tested bonds admit
one common neighborhood. Local dependence of $f$ then gives
$f(W(y))=f(\exp(Z(y))W_0)$ there. No compatibility between different
$g_i$ is needed, since $Z$ is an ordinary fine bond field in the original
coordinates. In particular this argument does not require a logarithm
of $W_0$ around a noncontractible loop.

Here is the bound on this field. Write
$\omega(i)=\theta_f(1+\operatorname{dist}(i_-,o))^{p_f}$ for the
fine direction width. Let $C\ge1$ be the fixed local gauge response
constant, $M\ge1$ the weighted tree-response constant just specified, and put
\[
 H_X=(2+\operatorname{breadth}(X))^{np_c},\qquad
 A_X=2MC+\frac{(c_nQ+M)s^2H_X}{\theta_f}.
\]
The first derivative of $b_i$ obeys the coarse weighted direction bound
with cost $C$; the linear response estimate and
$\theta_c\le2\theta_f$, $p_c\le p_f$, therefore give
$2MC\omega(i)$. For derivatives of order $j\ge2$, the local chart
bound gives
\[
 \|D^jb_i(0)\|\le s^j
 (1+\operatorname{dist}(\operatorname{parent}(i_-),
                         \operatorname{parent}(o)))^{jp_c}
 \le s^2H_X.
\]
The last inequality uses $s\le1$, $j\le n$, and the diameter bound
for a tested bond relative to $o\in X$. The linear response on this
bounded field costs at most $Ms^2H_X$.

The relative display on two copies of its inputs vanishes at zero;
its derivative there is the linear response of their difference, up to
$Q\theta_c$. Its first derivative away from zero differs by at most
a constant times $Qs$, and derivatives of higher order are bounded by
$Q$. The partition chain rule, exactly as for
\eqref{eq:reference-display-remainder}, bounds the remaining terms
at order $j\ge1$ by $c_nQs^2H_X$. Adjoint return preserves the norm.
Since $\omega(i)\ge\theta_f$, these estimates prove
\[
                  \|D^jZ_i(0)\|\le A_X\omega(i),
                         \qquad 1\le j\le n.
\]
The cutoff and the local health estimate make $W_0$ healthy. Apply the
ordinary fine regularity bound to $f(\exp(Z)W_0)$ at $Z=0$.
For each partition its argument fields have the displayed fine widths,
so that the bound is at most $C_n\tau A_X^n$ after summing partitions.
The zeroth derivative is bounded by $\tau$. This proves a general-support
bound whose leading constant $C_n(2MC)^n$ is chosen before $L$.
Every dependence on the later chart constant $Q$ is accompanied by
$s^2/\theta_f$, which can be made small after $L$ is fixed.
There is no sum of separate insertion derivatives over bonds: the
entire field $Z$ is inserted into the same regularity estimate.

\subsection{Integrating the density derivatives}
Use the same finite Gaussian coordinates for all derivatives,
with its positive covariance and the frozen fine field fixed.
The moment and density differentiation arguments of
Sections~\ref{sec:gaussian-differentiation} and
\ref{sec:gaussian-moment-prices} apply to this reference as well:
its precision floor, inverse jets and local affine source were
proved in Section~\ref{sec:principal-comparison-jets}. The cutoff
is held fixed in these coordinates. Only on its support is the
unmasked regular insertion evaluated; the health conclusion in
the preceding paragraph applies there.

We detail the dependence of the bound on the support. Let $q$
be its fine cover number and let
$H=(r_0+48)(1+B)$, where $r_0$ is the fixed allocation radius.
The number of read coordinates and all polynomial weights on
this set are bounded by fixed powers of $H(1+q)$, with constants
allowed to depend on $L$. Differentiating a Gaussian density
inserts sums of products of inverse-covariance entries and at
most quadratic coordinate factors. At order at most $n$ the
result is a polynomial in the Gaussian coordinates of degree
at most $2n$. The derivative bounds and the product rule therefore
bound the price in \eqref{eq:reference-short-pointwise}, including
the density factors, by
\begin{equation}\label{eq:reference-short-envelope}
 C_0\tau L^{-d-1}
        \{1+\rho^{2n}(1+\|z\|)^{4n}\}.
\end{equation}
Here $C_0$ depends on the pre-$L$ graded-field cost and the
regularity conversion constants, while
\[
 0\le\rho\le K_L(\theta_f+g/\theta_f)
                         \{H(1+q)\}^{a_n}.
\]
To check this factor, insert $w=gz$ into
\eqref{eq:reference-graded-cost}. The nonlinear remainder
$s^2/\theta_f$ is bounded by $\theta_f$ times a polynomial in
$B$ because the two thresholds have bounded ratio. The noise
contributes $(g/\theta_f)\|z\|$. Every positive-order density
or covariance variation has at least one coarse direction factor,
which is at most $\theta_c$ times such a polynomial; enlarge
$K_L$ and $a_n$ to include it. Separate the fixed term $M$ from
these errors, and use
$(a+b)^{2n}\le(2a)^{2n}(1+b^{2n})$ for $a\ge1$, $b\ge0$.
This gives \eqref{eq:reference-short-envelope}. For a general support,
use the separate-gauge bound just proved, with its leading constant
$C_n(2MC)^n$ and error $(c_nQ+M)s^2H_X/\theta_f$.
Combining it with the same density factors gives $C_1\tau$ in
place of $C_0\tau L^{-d-1}$, with possibly larger fixed exponents.
Its pre-$L$ constants are the ordinary local regularity and
linear-response constants.

The Gaussian moment of $(1+\|z\|)^{4n}$ is bounded by
$K'_L\{H(1+q)\}^{b_n}$. This follows by expanding the norm
bound into coordinate moments and using the covariance upper
bound, exactly as in Section~\ref{sec:gaussian-moment-prices}.
Choose $h_n=2na_n+b_n$. For each fixed $L$, increasing the
entrance coupling makes
\[
 [K_L(\theta_f+g/\theta_f)]^{2n}K'_L H^{h_n}\le1.
\]
Indeed $\theta_f$ times any fixed power of $B$ tends to zero,
and, uniformly over the incoming beta window,
$g/\theta_f\le2/\operatorname{monitor}(\beta_f)$.
The monitor grows faster than every fixed power of $B$.
This argument does not assign a negative power exponent to
$g/\theta_f$. The resulting integrated derivative bounds are
\begin{equation}\label{eq:reference-integrated-attenuation}
 \begin{cases}
 C\tau L^{-d-1}(1+q)^{h_n},&q^3<L,\\
 C'\tau(1+q)^{h'_n},&q^3\ge L.
 \end{cases}
\end{equation}
All differentiation under the integral is justified by the same
polynomial Gaussian majorants on a common parameter neighborhood.

\subsection{Absorbing the enlarged support}
Let $A>0$ be the fine support exponent, $e\ge0$ the requested
coarse exponent including its reserve, and $q'$ the coarse
cover number of the allocated output support. The geometric
cover estimate gives
\[
 Lq'\le(L+10q)t_{r_0},
\]
where $t_{r_0}$ depends only on the fixed allocation radius.
Choose $L$ so that $20e t_{r_0}\le LA$. Then
\begin{equation}\label{eq:reference-cover-release}
 eq'\le e t_{r_0}+\tfrac12 A q.
\end{equation}
For every fixed integer $h$, the function
$(1+q)^h e^{-Aq/2}$ is bounded on $q\ge0$: it tends to zero
at infinity and is continuous on compact intervals. Thus the
short-support case of \eqref{eq:reference-integrated-attenuation}
and \eqref{eq:reference-cover-release} yield
\[
 e^{eq'}\,\|\hbox{reference integral}\|_{J^n}
                  \le C_2L^{-d-1}\tau e^{Aq}.
\]
For the long-support case, $q^3\ge L$ implies
$1\le((1+q)^3/L)^{d+1}$. Insert this inequality into its
bound, then absorb the larger fixed polynomial in the same
half of the support exponential. We obtain the identical form
with another constant $C_3$, also chosen before $L$.
This is the required gain for all supports, including wrapping
ones in the long-support case.

Finally, the allotted regular price has the form
$L^{-d-2/5}/(2h)$, where $h>0$ depends only on the fixed
profile. Since $L^{-d-1}/L^{-d-2/5}=L^{-3/5}$, choose $L$
large enough that both $C_2L^{-d-1}$ and $C_3L^{-d-1}$ are
at most half that price. All nonlinear chart, moment and
Gaussian-comparison thresholds are chosen afterwards. This
preserves the order of constants needed for the one-step
regular estimate.
